% ============================================================
% 202_FFGFT_Feldtheorie_Gesamt_En_ch.tex
% The Complete Field Theory of FFGFT / T0 Theory
% A coherent presentation
% Johann Pascher — May 2026
% ============================================================

%------------------------------------------------------------
\part{Foundations and Motivation}
%------------------------------------------------------------

\section{What is FFGFT?}

FFGFT (Fundamental Fractal Geometric Field Theory), also called
T0 Time-Mass Duality, describes all known elementary particles
as bound oscillation modes of a \textbf{universal fractal torus}.
The central idea: observed masses and interactions follow directly
from geometry, without arbitrary parameters.

The theory requires only two fundamental inputs:

\begin{itemize}
  \item \textbf{The geometric constant}
    $\xi_0 = \frac{4}{3}\times 10^{-4}$ — not postulated but
    derived from the tetrahedral symmetry of the 4D torus and the
    self-consistency of the electron-proton-muon system (Doc.~180).
  \item \textbf{The Planck length} $\ell_P$ — the natural length scale
    of quantum gravity.
\end{itemize}

All physical quantities — particle masses, coupling constants,
cosmological parameters — follow from these two quantities.

\section{The Fundamental Scale of the Torus}

The characteristic length scale of the universal torus:
\begin{equation}
  L_0 = \frac{\ell_P}{\xi_0}
  \approx 7500 \cdot \ell_P
  \approx 6.14 \times 10^{-16}\,\text{GeV}^{-1}
\end{equation}

This is the \textbf{fundamental period} of the universal torus.
All particles are modes on this torus — like overtones of a vibrating
string, but in three spatial dimensions. The Compton wavelength
$L_i = 2\pi/m_i$ is the wavelength of the respective mode,
not the size of a separate torus.

\section{The Fractal Dimension}

An ordinary torus has dimension 3. The FFGFT torus has a
\textbf{fractal correction}:
\begin{equation}
  D_f = 3 - \xi_0 \approx 2.9998667
\end{equation}

This minimal deviation is the origin of all particle masses.
It generates a \textbf{non-linear dispersion relation} — analogous
to an anharmonic crystal lattice. On a simple torus ($D_f = 3$)
masses would be proportional to wave number: $m_n \propto n$,
ratios integer. In reality:
\[
  \frac{m_\mu}{m_e} \approx 206.8, \quad
  \frac{m_\tau}{m_e} \approx 3477
\]
These non-integer ratios require the fractal dispersion.

%------------------------------------------------------------
\part{The Torus Modes and Their Quantum Numbers}
%------------------------------------------------------------

\section{Quantum Numbers of the Modes}

Each mode on the torus is characterised by three quantum numbers:
$n$ (principal), $l$ (orbital angular momentum, transverse component)
and $j$ (spin, $j = l \pm \tfrac{1}{2}$ for fermions).

The wave vector of a mode is quantised:
\begin{equation}
  \mathbf{k}_{n,l} = \frac{2\pi}{L_0}(k_x, k_y, k_z),
  \quad k_x, k_y, k_z \in \mathbb{Z}
\end{equation}

Connection to the quantum numbers:
\begin{align}
  n &= k_x^2 + k_y^2 + k_z^2 \\
  &\quad\text{(principal = squared total length)} \nonumber \\
  l &= \sqrt{k_x^2 + k_y^2} \\
  &\quad\text{(orbital = transverse component)} \nonumber
\end{align}

The $r_i$ prefactors and $p_i$ exponents are discrete winding ratios
of the 4D torus — not continuous parameters. For the leptons:

\begin{center}
{\footnotesize
\begin{tabular}{lccccl}
\toprule
\textbf{Particle} & \textbf{n} & \textbf{l} &
  \textbf{$(k_x,k_y,k_z)$} & \textbf{$r_i$} & \textbf{$p_i$} \\
\midrule
Electron & 1 & 0 & $(1,0,0)$ & $4/3$ & $3/2$ \\
Muon     & 2 & 1 & $(1,1,0)$ & $16/5$ & $1$ \\
Tau      & 3 & 2 & $(1,1,1)$ & $25/9$ & $2/3$ \\
\bottomrule
\end{tabular}
}
\end{center}

These values are the complete enumeration of discrete torus modes —
analogous to the integers in hydrogen, where one does not ask for
a formula for $n=1,2,3$. The table (Doc.~006) is the theory.

\section{Explicit Mode Functions}

On the torus $T^3 \times \mathbb{R}$, the eigenfunctions of the
d'Alembertian $\Box = \partial_t^2 - \nabla^2$ are:
\begin{equation}
  \boxed{
    \phi_{n,l,j}(t,\mathbf{x})
    = \frac{1}{\sqrt{V}}\,
      e^{i\mathbf{k}_{n,l}\cdot\mathbf{x}} \cdot
      Y_{l,j}(\theta,\varphi) \cdot
      e^{-im_{n,l,j}\,t}
  }
\end{equation}

with $V = L_0^3$ (torus volume) and spherical harmonics $Y_{l,j}(\theta,\varphi)$.

These functions are \textbf{orthonormal}:
\begin{equation}
  \int_{T^3} \phi_{n,l,j}^*(\mathbf{x})\,\phi_{n',l',j'}(\mathbf{x})\,d^3x
  = \delta_{nn'}\,\delta_{ll'}\,\delta_{jj'}
\end{equation}
and satisfy the \textbf{Klein-Gordon equation}:
\begin{equation}
  \Box\,\phi_{n,l,j} = -m_{n,l,j}^2\,\phi_{n,l,j}
\end{equation}

The spin structure follows from the winding numbers of the torus
(Doc.~051): spin-$s$ corresponds to torus winding $(n_\theta, n_\phi)$
with $n_\phi/n_\theta = 2s$. For spin-$\tfrac{1}{2}$: $(1,1)$;
the 720° symmetry follows directly from the topology.

%------------------------------------------------------------
\part{The Non-Closure Theorem}
%------------------------------------------------------------

\section{Algebraic Structure of the Modes}

The particle masses follow the geometric mass formula:
\begin{equation}
  m_{n,l,j} = r_{n,l,j} \cdot \xi_0^{p_{n,l,j}} \cdot v,
  \qquad v = 246\,\text{GeV}
  \label{eq:massformula}
\end{equation}

The prefactors $r_i$ are rational numbers (labels of torus modes);
the exponents $p_i$ are rational scale parameters.

\textbf{Restriction:} Numerical substitution is permitted.
Symbolic cross-multiplication is forbidden.

\section{The Theorem}

\begin{tcolorbox}[colback=blue!5, colframe=blue!60!black, boxrule=0.8pt,
  arc=3pt, left=4mm, right=4mm,
  title={Non-Closure Theorem of FFGFT Torus Modes}]
For all pairs of different torus modes $i \neq j$:
\[
  \forall\, i \neq j:\;
  (r_i \cdot r_j,\; p_i + p_j) \notin \mathcal{P}
\]
where $\mathcal{P}$ is the set of all admissible pairs $(r,p)$ in FFGFT.
The set $\mathcal{P}$ is \textbf{not closed} under cross-operations
between different torus modes.
\end{tcolorbox}

\textbf{Proof of violation (example):}
\begin{align*}
  r_e \cdot r_\mu &= \frac{4}{3} \cdot \frac{16}{5} = \frac{64}{15}
  \approx 4.267 \\
  p_e + p_\mu &= \frac{3}{2} + 1 = \frac{5}{2}
\end{align*}

In the complete mode table (Doc.~006), no particle has $r = 64/15$
and $p = 5/2$. The pair $(64/15, 5/2) \notin \mathcal{P}$.

\textbf{Why does the theorem hold?} The cause lies in the orthogonality
of the modes. In Hilbert space, the product of two different basis
vectors is structurally meaningless. Likewise, $r_e \cdot r_\mu$ is
computable as a number but belongs to no torus mode.

\textbf{Practical rule:}

\begin{center}
{\footnotesize
\begin{tabular}{p{5cm}p{5cm}}
\toprule
\textbf{Permitted} & \textbf{Forbidden} \\
\midrule
Numerical: $m_e \cdot m_\mu \approx 53.9\,\text{MeV}^2$ &
  Symbolic: $r_e \cdot r_\mu = 64/15$ as mass formula \\[4pt]
$m_i = r_i \cdot \xi_0^{p_i} \cdot v$ for fixed $i$ &
  $(r_i \cdot r_j) \cdot \xi_0^{p_i+p_j} \cdot v$ for $i \neq j$ \\[4pt]
Mass ratios $m_i/m_j$ numerically &
  Cross-multiplication of different $r_i, r_j$ \\
\bottomrule
\end{tabular}
}
\end{center}

%------------------------------------------------------------
\part{The Lagrangian and the Bridge to the Mode Operator}
%------------------------------------------------------------

\section{The Free Lagrangian}

Starting point is the free Klein-Gordon Lagrangian:
\begin{equation}
  \mathcal{L}_0 = \frac{1}{2}(\partial\delta m)^2
\end{equation}
The scalar field $\delta m(x)$ describes mass fluctuations on the torus.
The equation of motion $\Box\,\delta m = 0$ describes massless oscillations.

\begin{tcolorbox}[colback=gray!5, colframe=gray!50!black, boxrule=0.5pt,
  arc=2pt, left=3mm, right=3mm, fontupper=\small,
  title={Related document: Doc.~201 (DVFT adaptation)}]
Doc.~201 (\textit{FFGFT-alles}) develops the same starting point
$\mathcal{L}_0 = \varepsilon(\partial\Delta m)^2$ in a different
direction: the polar decomposition $\Phi = \rho e^{i\theta}$ with
vacuum amplitude $\rho$ and vacuum phase $\theta$ (DVFT ---
Dynamic Vacuum Field Theory). Applications: \textit{apparent} cosmic acceleration
(without dark energy), galactic rotation curves
(without dark matter), CMB homogeneity. No contradiction with Doc.~202 ---
different development directions of the same foundation.

\medskip

\textbf{Quantum-mechanical equations of motion:} The DVFT line in
Doc.~201 \S 2.6 (simplified Dirac equation from T0) and \S 6.6
(Schrödinger equation derivation from T0) provides the bridge to
non-relativistic and relativistic quantum mechanics. In the mode
picture of Doc.~202, the DVFT polar decomposition corresponds to the
low-energy limit of the mode superposition; the explicit bridge
formula is documented in the section
\enquote{Connection to the Quantum Mechanical Equations of Motion}
below.

\medskip

\textbf{Note on the terminology \enquote{renormalisation group}:}
The scale running of the coupling in Doc.~202 \S\,18 is referred to
there -- following standard QFT language -- as the
\emph{renormalisation group}. In FFGFT, however, this scale
structure follows directly from the recursion $\mathcal{R}$
(Doc.~203) on the fractal torus geometry and is \emph{not} the
result of a renormalisation procedure. Doc.~201 consequently
avoids this terminology and describes the scale flow directly
through the DVFT phase dynamics $(\rho, \theta)$ -- i.e.\ in an
FFGFT-native manner. The unified language convention (scale
structure from recursion, $\mathcal{R}$-induced scale correction
instead of RG language) is recorded in Doc.~190~R7 and is binding
for the entire series.
\end{tcolorbox}

\section{The Torus Correction Term}

The fractal geometry adds a correction term generating the non-linear dispersion:
\begin{equation}
  \Delta\mathcal{L}_{\text{Torus}}
  = \varepsilon\,\xi_0\!\left[
    \frac{f(\theta)}{2}(\partial\delta m)^2
    - k^{\mu\nu}(\theta)\,\partial_\mu\delta m\,\partial_\nu\delta m
    \right]
\end{equation}

The geometric functions $f(\theta)$ and $k^{\mu\nu}(\theta)$ encode
the curvature of the fractal torus. The complete FFGFT Lagrangian:
\begin{equation}
  \boxed{
    \mathcal{L} = \mathcal{L}_0 + \Delta\mathcal{L}_{\text{Torus}}
  }
\end{equation}

\section{The Mode Operator}

The Hilbert space of torus modes:
\begin{equation}
  \mathcal{H} = \bigoplus_{n,l,j} \mathbb{C}\cdot\mathbf{e}_{n,l,j},
  \qquad
  \langle\mathbf{e}_{n,l,j},\mathbf{e}_{n',l',j'}\rangle
  = \delta_{nn'}\delta_{ll'}\delta_{jj'}
\end{equation}

The field operator:
\begin{equation}
  \Phi(x) = \sum_{n,l,j} m_{n,l,j}\cdot\phi_{n,l,j}(x)\cdot P_{n,l,j}
\end{equation}
with $P_{n,l,j} = |\mathbf{e}_{n,l,j}\rangle\langle\mathbf{e}_{n,l,j}|$.
Since $P_iP_j = \delta_{ij}P_i$: $\Phi^2 = \sum_i m_i^2 P_i$.

\section{The Bridge Formula}

The Lagrangian and mode operator describe the same physics:
\begin{equation}
  \boxed{
    \int\frac{1}{2}(\partial\delta m)^2\,d^4x
    = \frac{1}{2}\langle\psi|\Phi^2|\psi\rangle
  }
  \label{eq:bridge}
\end{equation}

\textbf{Derivation in 6 steps:}

\textit{Step 1 — Mode expansion:} $\delta m(x) = \sum_i a_i\phi_i(x)$.

\textit{Step 2 — Substitution:}
$\int\mathcal{L}_0\,d^4x
= \tfrac{1}{2}\sum_{i,j}a_ia_j\int(\partial\phi_i)(\partial\phi_j)\,d^4x$.

\textit{Step 3 — Orthogonality:}
Cross terms vanish for $i\neq j$, leaving
$\tfrac{1}{2}\sum_i a_i^2\int(\partial\phi_i)^2\,d^4x$.

\textit{Step 4 — Klein-Gordon:}
$\int(\partial\phi_i)^2\,d^4x = m_i^2$ (partial integration, normalisation).

\textit{Step 5 — Summary:}
$\int\mathcal{L}_0\,d^4x = \tfrac{1}{2}\sum_i a_i^2 m_i^2$.

\textit{Step 6 — Mode operator:}
$\langle\psi|\Phi^2|\psi\rangle = \sum_i a_i^2 m_i^2$.
Comparison gives \eqref{eq:bridge} with $\varepsilon = 1/2$.

%------------------------------------------------------------
\part{Mass Generation as a Binding Phenomenon}
%------------------------------------------------------------

\section{The Binding Picture}

The eigenvalue equation on the torus:
\begin{equation}
  \left[\Box + \hat{V}(\xi_0, D_f)\right]\phi_{n,l,j} = -m_{n,l,j}^2\,\phi_{n,l,j}
\end{equation}

The mass is a difference:
\begin{equation}
  m_i^2
  = \underbrace{\left(\frac{2\pi n}{L_0}\right)^2}_{\text{kinetic}}
  + \underbrace{\langle\phi_i|\hat{V}|\phi_i\rangle}_{\approx\,-(2\pi n/L_0)^2}
\end{equation}

The kinetic energy is enormous ($\sim 10^{32}\,\text{GeV}^2$), the potential
almost exactly opposite. The mass is the tiny residual:

\begin{center}
{\footnotesize
\begin{tabular}{lccc}
\toprule
\textbf{Particle} & \textbf{$(2\pi n/L_0)^2$} &
  \textbf{$m_i^2/\text{(kin.)}$} \\
\midrule
Electron & $1.046\times10^{32}\,\text{GeV}^2$ & $2.5\times10^{-39}$ \\
Muon     & $4.185\times10^{32}\,\text{GeV}^2$ & $2.7\times10^{-35}$ \\
Tau      & $9.415\times10^{32}\,\text{GeV}^2$ & $3.4\times10^{-33}$ \\
\bottomrule
\end{tabular}
}
\end{center}

\section{Not Fine-Tuning --- a Binding Phenomenon}

The extreme cancellation (32–39 orders of magnitude) is not fine-tuning
— it is the defining property of bound states.

\textbf{Crystal analogy:} One does not ask why a crystal produces a
particular phonon dispersion. The dispersion curve \emph{is} the crystal
structure in momentum space.

\textbf{Hydrogen analogy:} Kinetic energy $\approx 13.6\,\text{eV}$,
Coulomb potential $\approx -27.2\,\text{eV}$, binding energy $-13.6\,\text{eV}$.
Nobody asks why the potential compensates the kinetic energy so precisely —
it follows from the Schrödinger equation.

\textbf{In FFGFT:} The particles are bound torus modes. The mass is
the binding residual. The potential $\hat{V}$ is uniquely determined by
the torus geometry ($\xi_0$, $D_f$) — no free parameter.

\section{Connection to the Lagrangian Extension}

\begin{center}
{\scriptsize
\resizebox{\textwidth}{!}{\begin{tabular}{p{5cm}p{5cm}}
\toprule
\textbf{Lagrangian formulation} & \textbf{Torus binding picture} \\
\midrule
$\xi_0 \cdot m_\ell \cdot \bar\psi\psi \cdot \delta m$ (coupling) &
  Coupling of fermion mode to its own torus mode \\[4pt]
$\Delta\mathcal{L}_{\text{Torus}}$ &
  Potential term $\hat{V}$ from $D_f = 3-\xi_0$ \\[4pt]
Mass parameter $m_\ell$ appears explicitly &
  Mass emerges as eigenvalue residual \\[4pt]
Standard + geometric correction &
  Kinetic $+$ potential $\approx 0$, residual $= m_i^2$ \\
\bottomrule
\end{tabular}}}
\end{center}

\textbf{Difference from the Higgs mechanism:}
In the Standard Model, the Yukawa coupling $y_i$ is a free parameter.
In FFGFT, the corresponding term $\xi_0 \cdot m_\ell$ is geometrically
fixed by $\xi_0 = 4/30000$.

The central claim:
\begin{quote}
\emph{Standard QFT calculations must be extended by a tiny geometric
mass contribution that arises directly from the torus structure.}
\end{quote}

\section{Connection to the Quantum Mechanical Equations of Motion}

The mode structure developed in Doc.~202 and the free Lagrangian
$\mathcal{L}_0 = \tfrac12(\partial\delta m)^2$ provide the
\emph{spectral framework}: masses arise as eigenvalues of the mode
operator, not as Lagrangian parameters. For the quantum-mechanical
equations of motion the parallel DVFT line in Doc.~201 is invoked,
which develops the same starting point
$\mathcal{L}_0 = \varepsilon(\partial\Delta m)^2$ in a different
direction.

\subsection*{Bridge 202 $\leftrightarrow$ 201}

In the low-energy limit, the mode sum of Doc.~202 maps to the DVFT
polar decomposition of Doc.~201:
\begin{equation}
  \Phi(x) \;=\; \sum_{n,l,j} m_{n,l,j}\,\phi_{n,l,j}(x)\,P_{n,l,j}
  \;\;\xrightarrow{\;\text{low-energy}\;}\;\;
  \sqrt{\rho(x,t)}\,e^{i\theta(x,t)}
  \tag{202.\,QM-1}
\end{equation}
with
$\rho(x,t) = \big|\sum_{n,l,j} \alpha_{n,l,j}\phi_{n,l,j}(x)\big|^2$
as energy density and $\theta(x,t)$ as the coherent phase.

\textit{Comprehensive sources for the QM extension:} Doc.~020
(T0-QM-QFT-RT, comprehensive treatment), Doc.~035 (QM overview).

\subsection*{Simplified Dirac equation (from Doc.~201 \S\,2.6)}

On this basis, Doc.~201 replaces the full Dirac equation by:
\begin{equation}
  \partial^{2}\,\Delta m \;=\; 0
  \tag{202.\,QM-2}
\end{equation}
The fundamental structure here is the Clifford algebra
$\mathbf{e}_\mu\mathbf{e}_\nu + \mathbf{e}_\nu\mathbf{e}_\mu = 2g_{\mu\nu}$
(Doc.~050, 051) rather than a specific matrix representation. The
$4\times 4$ $\gamma$ matrices nevertheless remain \emph{as a
representation} of the Clifford generators and stay a practical tool
for concrete calculations --- they are therefore not eliminated but
rather traced back to their geometric meaning. Spin manifests itself
on the fractal torus as a topological property
(\emph{vacuum phase winding}), consistent with the quantum number
$j = l \pm \tfrac12$ from the mode catalogue of Doc.~202 \S
\enquote{Quantum Numbers of the Modes}. A relativistic state becomes
\begin{equation}
  \Phi(x,t) \;=\; \rho(x,t)\,e^{i\theta(x,t)}\cdot \chi(\sigma)
  \tag{202.\,QM-3}
\end{equation}
with $\chi(\sigma)$ as the spin component arising from geometric
knot modes.

\textit{Parallel derivations on the Dirac line:}
Doc.~020 \S\,\enquote{T0-modified Dirac equation},
Doc.~067 \S\,\enquote{Modified Dirac equation},
Doc.~095 \S\,\enquote{Simplified Dirac equation (T0 version)},
Doc.~053 (mass elimination in the Dirac Lagrangian).

\subsection*{Schrödinger equation (from Doc.~201 \S\,6.6)}

In the non-relativistic limit, the phase equation of the vacuum field
reduces to the standard Schrödinger equation:
\begin{equation}
  i\hbar\,\frac{\partial\psi}{\partial t}
  \;=\; -\frac{\hbar^{2}}{2m}\,\nabla^{2}\psi + V\psi
  \tag{202.\,QM-4}
\end{equation}
with $\psi \propto \sqrt{\rho}\,e^{i\theta}$, effective mass $m$ from
the knot inertia, and potential $V$ from external $\rho$ perturbations.
The Schrödinger equation is therefore \emph{not fundamental} but an
effective equation for slow, non-relativistic knot excitations.

\textit{Parallel derivations on the Schrödinger line:}
Doc.~020 \S\,\enquote{T0-modified Schrödinger equation},
Doc.~067 \S\,\enquote{Modified Schrödinger equation},
Doc.~095 \S\,\enquote{Extended Schrödinger equation (T0-modified)},
Doc.~129 \S\,\enquote{Schrödinger equation in simplified T0 form},
Doc.~037 (Cairo counterexample and geometric resolution).

\subsection*{Bell correlations (from Doc.~023, 074)}

Entangled states correspond in FFGFT to correlated mode configurations
sharing a common phase winding on the torus. The modified CHSH form
contains geometric constants that follow directly from the Lagrangian
framework:
\begin{equation}
  \mathrm{CHSH}^{\mathrm{FFGFT}}
  \;=\; 2\sqrt{2}\cdot K_{\text{frak}}^{D_f}\cdot
        \left(1 - \xi_0\cdot\tfrac{\Delta\theta}{\pi}\right)
  \tag{202.\,QM-5}
\end{equation}
with $D_f = 3 - \xi_0$ (fractal dimension from the torus correction
term $\Delta\mathcal{L}_{\text{Torus}}$, \S\,9), $K_{\text{frak}}$ as
geometric correction factor, and $\xi_0 = 4/30000$. The prediction is
a tiny deviation from the standard quantum value:
\begin{equation}
  \Delta\mathrm{CHSH} \;=\; 2\sqrt{2} - \mathrm{CHSH}^{\mathrm{FFGFT}}
  \;\sim\; \mathcal{O}(\xi_0)\;\approx\; 10^{-4}
  \tag{202.\,QM-6}
\end{equation}

\textbf{Lagrangian connection.} Bell correlations do not follow
directly from $\mathcal{L} = \mathcal{L}_0 + \Delta\mathcal{L}_{\text{Torus}}$,
but from the \emph{topology} of the mode configurations. The
Lagrangian nevertheless supplies the constants
$\xi_0$, $D_f$ and $K_{\text{frak}}$ appearing in (202.\,QM-5).
Bell tests thereby become an \emph{indirect} probe of the Lagrangian
structure via the $\xi_0$ scale constant.

\textbf{Locality picture.} The correlations arise from topological
knot correlation on the torus, not from non-local action. Entanglement
appears as a shared phase winding between two mode sectors with
correlated quantum numbers $(n,l,j)$. Experimentally accessible in
loophole-free tests at large angles $\Delta\theta > \pi$ with
resolution $\sim 10^{-4}$ (Doc.~023, 074, 131).

\textit{Parallel derivations on the Bell line:}
\begin{itemize}\setlength\itemsep{0pt}
  \item Dedicated Bell series:
        Doc.~023 (Bell tests in the T0 context, modified CHSH form),
        Doc.~023a Bell-Part 2 (multi-qubit extension, philosophical
        reflections, experimental proposals),
        Doc.~023a Bell-Video (locally realistic picture in detail).
  \item Quantum computing application:
        Doc.~147 (Bell test modifications \S, 73-qubit predictions
        $\Delta\mathrm{CHSH} \sim 10^{-3}$, Bell-enhanced peak
        detection).
  \item Locality \& non-locality:
        Doc.~074 (no-go theorem and T0 answer to Bell),
        Doc.~131 (\enquote{apparently instantaneous} -- locality picture),
        Doc.~055 (DynMassePhotons non-local).
  \item Cross-cutting:
        Doc.~022 (T0-QFT-ML addendum),
        Doc.~035 \S\,\enquote{Bell tests \& EPR},
        Doc.~175 (qubit state spaces).
\end{itemize}

\subsection*{Qubit mappings (from Doc.~175, 034)}

A qubit state is, in the FFGFT picture, a specific configuration of
the mode structure. From the bridge formula (202.\,QM-1) one obtains
two localised basis configurations:
\begin{equation}
  \Phi_0(x) \;=\; \rho_0(x)\,e^{i\theta_0(x,t)},
  \qquad
  \Phi_1(x) \;=\; \rho_1(x)\,e^{i\theta_1(x,t)}.
  \tag{202.\,QM-7}
\end{equation}
The coherent superposition
$\alpha\Phi_0 + \beta\Phi_1$ ($|\alpha|^2 + |\beta|^2 = 1$) yields the
general qubit state as a \emph{single phase configuration} of the
vacuum field --- not an ontological multiple existence but a coherent
superposition of modes.

\textbf{Cylindrical phase space.} The natural geometry of a qubit in
FFGFT (Doc.~034, 175) is not the Bloch sphere but a cylindrical
phase space with three coordinates $(z, r, \theta)$:
\begin{itemize}\setlength\itemsep{0pt}
  \item $z \in [-1, +1]$: torsion configuration along
        $\ket{0} \leftrightarrow \ket{1}$
  \item $r \in [0, 1]$: strength of intermediate torsion
  \item $\theta \in [0, 2\pi)$: relative phase
\end{itemize}
with the Bloch-sphere correspondence $z = \cos\theta_{\text{Bloch}}$,
$r = |\sin\theta_{\text{Bloch}}|$. The consequence: \emph{quantum
gates become geometric transformations in real space} --- the Hadamard
gate exchanges $z \leftrightarrow r$ with a $\pi/2$ phase rotation;
the Z phase is a pure rotation about the cylinder axis.

\textbf{Lagrangian connection.} The basis states $\ket{0}, \ket{1}$
are two specific mode excitations $\phi_{n,l,j}$ from the mode
catalogue of Doc.~202 \S\,\enquote{Quantum Numbers of the Modes}
and \S\,\enquote{Explicit Mode Functions}. The geometry of the mode
functions arises from the Lagrangian
$\mathcal{L} = \mathcal{L}_0 + \Delta\mathcal{L}_{\text{Torus}}$;
qubit operations are effective transformations on the Hilbert space
spanned by the mode structure. Spin qubits sit at the $\xi$ level
$N \approx 8$ (nanometre scale, quantum confinement); photonic
resonators sit on a different level of the $\xi$ hierarchy
(Doc.~175).

\textit{Parallel derivations on the qubit line:}
\begin{itemize}\setlength\itemsep{0pt}
  \item Geometry and phase space:
        Doc.~034 (QM optimisation, introduction of the cylindrical
        phase space),
        Doc.~175 (qubit state spaces, quantum registers),
        Doc.~173 (quantum-dot resonator, level $N \approx 8$).
  \item Spin manipulation and readout:
        Doc.~174 (five methods for spin manipulation, QND
        measurement),
        Doc.~178 (FFGFT spin stability).
  \item Quantum computing architecture:
        Doc.~147 ($\phi$-hierarchical QFT, Shor's algorithm,
        deterministic energy-field qubits),
        Doc.~176 (Shor-photon),
        Doc.~161 (Ising machine),
        Doc.~183 (Willow self-measurement).
\end{itemize}

\subsection*{Consequence for the Verification Condition}

Both equations of motion (202.\,QM-2) and (202.\,QM-4) are contained
in the mode picture of Doc.~202 -- as different low-energy limits of
the same spectral structure. The verification condition of the next
section therefore applies consistently to both sectors. The Bell
correlation (202.\,QM-5) and the qubit mapping (202.\,QM-7) provide
independent consistency tests of the same $\xi_0$ constant and the
same mode structure.

\section{The Verification Condition}

The theory makes a numerically testable prediction:
\begin{equation}
  \boxed{
    \langle\phi_{n,l,j}|\hat{V}|\phi_{n,l,j}\rangle
    = (r_i)^2 \cdot \xi_0^{2p_i} \cdot v^2
      - \left(\frac{2\pi n}{L_0}\right)^2
  }
  \label{eq:verif}
\end{equation}

This is the Stage~3 bottleneck (Doc.~189): explicitly derive
$\Delta\mathcal{L}_{\text{Torus}}$ from the torus geometry and show
that the eigenvalues match the mass table (Doc.~006).

\section{Methodological Status of the Verification Condition}

The verification condition (\ref{eq:verif}) is evaluated against
experimental masses. This evaluation carries three methodological
limitations that must be made explicit, since they are often passed
over.

\subsection*{Reduction is not precision improvement}

The FFGFT Lagrangian
$\mathcal{L} = \mathcal{L}_0 + \Delta\mathcal{L}_{\text{Torus}} + \mathcal{L}_{\text{Yukawa}}$
contains the Standard-Model Lagrangian as a limiting case: for
$\xi_0 \to 0$ both $\Delta\mathcal{L}_{\text{Torus}}$ and the Yukawa
term $\xi_0\cdot m_\ell\cdot\bar\psi\psi\cdot\delta m$ vanish. Where
standard calculations are experimentally confirmed (e.g.\ QED
scattering amplitudes to $\sim 10^{-12}$), FFGFT can only reproduce
these and add $\xi_0$ corrections of order $10^{-4}$, not surpass them
qualitatively.

The FFGFT mass formulas $m_\ell = r_\ell\cdot\xi_0^{p_\ell}\cdot v$
yield values agreeing with PDG values to $\sim 1\,\%$ (Doc.~190~C5).
This is \emph{not} a precision result but a \emph{structural proof}:
the masses follow from the geometry, not from free parameters. The
reduction of $\sim 25$ free Standard-Model parameters to a single
geometric parameter $\xi_0$ is the quantifiable claim --- not the
numerical agreement of the individual masses.

\subsection*{Predictions beyond the Standard Model}

Where FFGFT delivers \emph{different} statements, these are
predictions for phenomena where the Standard Model is either silent
or predicts corrections not yet experimentally resolved:

\begin{itemize}\setlength\itemsep{0pt}
  \item Bell CHSH correction $\Delta\mathrm{CHSH} \sim \mathcal{O}(\xi_0)
        \approx 10^{-4}$ in loophole-free tests at large angles
        (Doc.~023, 074, 147; see \S\,15).
  \item $\Delta$CHSH$\sim 10^{-3}$ in 73-qubit systems with
        Bell-enhanced peak detection (Doc.~147).
  \item Modified dispersion at near-Planck energies (Doc.~055, 074).
  \item Dark matter and dark energy without separate components
        (DVFT line, Doc.~201).
  \item CMB homogeneity without inflation (Doc.~140, 201).
  \item Hierarchy problem: Higgs mass vs.\ Planck mass as a geometric
        consequence, not a fine-tuning problem (Doc.~003, 049).
\end{itemize}

\subsection*{Circularity of the comparison basis}

A systematic limitation of any numerical competition with the
Standard Model concerns the experimental comparison basis itself:
PDG values are not model-free raw data but already contain assumptions
about
\begin{itemize}\setlength\itemsep{0pt}
  \item units and gauge conventions,
  \item calibrating natural constants ($c$, $\hbar$, $e$ were fixed
        by definition in the 2019 SI reform),
  \item loop corrections from QFT evaluations.
\end{itemize}
These points are systematically treated in Doc.~066 (parameter-transfer
problem between unit systems) and Doc.~101 (circularity of the
constants).

An FFGFT calculation with a simpler entry point --- starting directly
from $\xi_0$ in natural units ($\hbar = c = 1$) --- avoids the
accumulation of model assumptions along the standard calibration
chain. As soon as it measures its values against PDG data, however,
it inherits their precorrections. The agreement at the $\sim 1\,\%$
level should therefore be read in two directions: as confirmation of
the geometric structural claim \emph{and} as a pointer to the
latitude introduced by the circular measurement basis. A direct
numerical competition would be fair only if the entire measurement
chain were conceptually reorganised --- a programme outlined in
Doc.~066 and 101 but exceeding the scope of the present field-theoretic
formulation.

\subsection*{Consequence for the assessment}

The verification condition (\ref{eq:verif}) is therefore to be read
\emph{structurally}, not \emph{quantitatively}: if the eigenvalues
of the mode-operator equation are confirmed to match the mass table
within the $\sim 1\,\%$ band, the geometric derivation is supported.
The remaining margin is not to be interpreted as a theoretical deficit
but as a reflex of the three methodological limitations stated above.

%------------------------------------------------------------
\part{The Field-Theoretic Formulation}
%------------------------------------------------------------

\section{Hilbert Space and Recursion Operator}

The recursion operator $\mathcal{R}$ formalises the fractal scaling structure:
\begin{equation}
  \mathcal{R}(\mathbf{e}_{n,l,j})
  = \xi_0^{p(n,l,j)-1} \cdot K_{\text{frak}} \cdot \mathbf{e}_{n,l,j},
  \qquad \mathcal{R}(\Phi_*) = \Phi_*
\end{equation}

The time-mass duality is a field symmetry:
$\mathcal{R}:\; t \to \xi_0 t,\; m \to \xi_0^{-1} m,\; \Phi \to \Phi$.

\section{Algebraic Structure}

Mode algebra for free waves on the torus:
\begin{equation}
  \mathbf{e}_{\mathbf{k}_i} \star \mathbf{e}_{\mathbf{k}_j}
  = \mathbf{e}_{\mathbf{k}_i + \mathbf{k}_j},
  \qquad
  C_{ij}^k = \delta_{k,\,i+j} \pmod{\text{torus period}}
\end{equation}

\section{Scale Structure from Recursion}

The recursive $\alpha$ running in FFGFT arises from the iterated
application of the recursion operator $\mathcal{R}$ (Doc.~203) to the
mode structure. In the formal writing of a modified QED $\beta$-function
(Doc.~008):
\begin{equation}
  \frac{d\alpha}{d\ln\mu}
  = \beta_0\alpha^2\cdot\left(1 + \xi_0\cdot\ln\frac{\mu}{E_0}\right)
\end{equation}
\begin{equation}
  \alpha^{-1}(\mu)
  = \alpha^{-1}(m_e)
    - \frac{\beta_0}{2\pi}\ln\frac{\mu}{m_e}
    - \frac{\beta_0\xi_0}{4\pi}\left(\ln\frac{\mu}{m_e}\right)^2
\end{equation}

The $\xi_0$ term is the FFGFT-specific addition. $\xi_0$ is \textbf{not}
a fixed point of this $\beta$-function but a fundamental geometric constant
--- time-invariant because topologically fixed.

\medskip

\textbf{Note on terminology.}
The scale running was occasionally referred to as the
\enquote{renormalisation group} in earlier FFGFT documents. This
designation is inaccurate: the scale flow here is \emph{not} the result
of a renormalisation procedure but a direct consequence of the recursion
$\mathcal{R}$ on the fractal torus geometry. Stability, scale flow and
UV finiteness are three aspects of the same recursive structure.
The binding language convention is recorded in Doc.~190~R7; the
FFGFT-native view is already developed in Doc.~159 (\enquote{without
artificial renormalisation}) and Doc.~189 (\enquote{geometrically
determined, not by renormalisation}).

%------------------------------------------------------------
\part{Open Mathematical Tasks}
%------------------------------------------------------------

\section{Four Precisely Formulated Problems}

The theory is conceptually closed. The remaining tasks are mathematical.

Problems 1 and 2 have been formally worked out in Doc.~203
($\mathcal{R}$-operator, fixed-point condition) and Doc.~204
(fractal boundary condition, first-order perturbation in $\xi_0$).
Problem 3 (algebra) is complete for the free theory.
Problem 4 (interactions) remains open.

\begin{center}
{\setlength{\hfuzz}{5pt}\scriptsize
\resizebox{\textwidth}{!}{\begin{tabular}{p{0.3cm}p{2.5cm}p{4.5cm}p{3.5cm}}
\toprule
\textbf{No.} & \textbf{Problem} & \textbf{Task} & \textbf{Status} \\
\midrule
1 & $\mathcal{R}$-operator &
  $\lambda_{n,l,j} = \xi_0^{p-1}K_{\text{frak}}$; $p(n,l,j)$ from torus &
  Doc.~203 \\[4pt]
2 & Mode functions &
  Incorporate fractal boundary condition &
  Doc.~204 \\[4pt]
3 & Algebra &
  Hopf algebra for interactions &
  Free theory: complete \\[4pt]
4 & Interactions &
  $\mathcal{L}_{\text{int}}$, $\beta$-function &
  Doc.~008; $\xi_0$ geometric \\
\bottomrule
\end{tabular}}}
\end{center}

%------------------------------------------------------------
\part{Summary}
%------------------------------------------------------------

\section{Overview of Results}

FFGFT / T0 theory describes all known elementary particles as modes
of a universal fractal torus — closed, consistent, from a single
geometric constant $\xi_0 = 4/3 \times 10^{-4}$.

\begin{center}
{\footnotesize
\resizebox{\textwidth}{!}{\begin{tabular}{p{3.8cm}p{7cm}}
\toprule
\textbf{Concept} & \textbf{Meaning} \\
\midrule
Universal torus &
  $D_f = 3-\xi_0$, fundamental scale $L_0 = \ell_P/\xi_0$ \\[4pt]
Torus modes &
  Quantum numbers $(n,l,j)$, wave vector
  $\mathbf{k} = (2\pi/L_0)(k_x,k_y,k_z)$ \\[4pt]
Mode functions &
  $\phi = V^{-1/2}e^{i\mathbf{k}\cdot\mathbf{x}}Y_{l,j}e^{-imt}$ \\[4pt]
Lagrangian &
  $\mathcal{L} = \frac{1}{2}(\partial\delta m)^2 +
  \Delta\mathcal{L}_{\text{Torus}}$ \\[4pt]
Mode operator &
  $\Phi = \sum_i m_i P_i$, diagonal in Hilbert space \\[4pt]
Bridge &
  $\int\mathcal{L}_0\,d^4x = \frac{1}{2}\langle\psi|\Phi^2|\psi\rangle$ \\[4pt]
Mass formula &
  $m_i = r_i\cdot\xi_0^{p_i}\cdot v$, $r_i,p_i \in \mathbb{Q}$ \\[4pt]
Non-closure &
  $(r_ir_j, p_i+p_j)\notin\mathcal{P}$ for $i\neq j$ \\[4pt]
Mass generation &
  Bound states: kin.~$\approx$ pot., mass $=$ binding residual \\
\bottomrule
\end{tabular}}}
\end{center}


%------------------------------------------------------------
\section*{Notation}
%------------------------------------------------------------

\begin{center}
{\footnotesize
\begin{tabular}{lp{8.5cm}}
\toprule
\textbf{Symbol} & \textbf{Meaning} \\
\midrule
$\xi_0 = 4/30000$ & Geometric base constant of the fractal torus \\
$\ell_P$ & Planck length \\
$L_0 = \ell_P/\xi_0$ & Fundamental torus scale \\
$D_f = 3-\xi_0$ & Fractal dimension of the FFGFT torus \\
$v = 246\,\text{GeV}$ & Higgs vacuum expectation value (energy unit) \\
$n,\,l,\,j$ & Quantum numbers: principal, orbital, spin \\
$\mathbf{k}_{n,l}$ & Quantised wave vector on the torus \\
$Y_{l,j}(\theta,\varphi)$ & Spherical harmonics \\
$\phi_{n,l,j}(t,\mathbf{x})$ & Torus mode functions \\
$m_i = r_i\cdot\xi_0^{p_i}\cdot v$ & FFGFT mass formula \\
$r_i \in \mathbb{Q}$ & Rational prefactor (winding ratio) \\
$p_i \in \mathbb{Q}$ & Rational exponent (scaling depth) \\
$\mathcal{P}$ & Set of admissible pairs $(r_i, p_i)$ \\
$\mathcal{H}$ & Hilbert space of torus modes \\
$\mathbf{e}_{n,l,j}$ & Basis vector of the Hilbert space \\
$P_i = |\mathbf{e}_i\rangle\langle\mathbf{e}_i|$ & Projector onto mode $i$ \\
$\Phi = \sum_i m_i P_i$ & Mode operator (field operator) \\
$\mathcal{L}_0 = \frac{1}{2}(\partial\delta m)^2$ & Free Lagrangian \\
$\Delta\mathcal{L}_{\text{Torus}}$ & Torus-geometric correction term \\
$\hat{V}(\xi_0, D_f)$ & Potential operator from torus geometry \\
$K_{\text{frak}} \approx 0.986$ & Fractal correction constant \\
$C^k_{ij} = \delta_{k,i+j}$ & Structure constants of the mode algebra \\
$\beta_0$ & Coefficient of the QED beta function \\
\bottomrule
\end{tabular}
}
\end{center}

\medskip\noindent
\textbf{References:} Doc.~006 (quantum numbers), Doc.~008 ($\beta$-function),
Doc.~049 ($L_0$), Doc.~051 (spinor), Doc.~129, Doc.~180 ($\xi_0$),
Doc.~189 (stage structure).
